\subsection{EXTENSION OF DERIVATIONS}

PROPOSITION 14. Let A be a commutative ring, M an A-module, B the A-algebra \(\mathsf{T}(\mathbf{M})\) (resp. \(\mathsf{S}(\mathbf{M})\) , resp. \(\bigwedge (\mathbf{M}))\) and E a (B, B)-bimodule. Let \(d_0: \mathbf{A} \to \mathbf{E}\) be a derivation of the ring A into the A-module E and \(d_{1}:M\to E\) an additive group homomorphism such that, for all \(a\in A\) and all \(x\in M\) ,

\[
d _ {1} (a x) = a d _ {1} (x) + d _ {0} (a). x,\tag{29}
\]

and further, when \(\mathbf{B} = \mathbf{S}(\mathbf{M})\) ,

\[
x. d _ {1} (y) + d _ {1} (x). y = y. d _ {1} (x) + d _ {1} (y). x\tag{30}
\]

for all \(x, y\) in \(\mathbf{M}\) , and, when \(\mathbf{B} = \bigwedge (\mathbf{M})\) ,

\[
x. d _ {1} (x) + d _ {1} (x). x = 0\tag{31}
\]

for all \(x \in \mathbf{M}\) . Then there exists one and only one derivation \(d\) of \(\mathbf{B}\) (considered as a \(\mathbf{Z}\) -algebra) into the \((\mathbf{B}, \mathbf{B})\) -bimodule \(\mathbf{E}\) such that \(d|_{\mathbf{A}} = d_0\) and \(d|M = d_1\) .

We take on the \(\mathbf{Z}\) -module \(\mathbf{B} \oplus \mathbf{E}\) the associative \(\mathbf{Z}\) -algebra structure defined by

\[
(b, t) (b ^ {\prime}, t ^ {\prime}) = (b b ^ {\prime}, b t ^ {\prime} + t b ^ {\prime})
\]

which has (1, 0) as unit element (no. 8, Proposition 13). Under the canonical injection \(t \mapsto (0, t)\) , E is identified with a two-sided ideal of B ⊕ E such that \(\mathrm{E}^2 = \{0\}\) . On the other hand, the mapping \(h_0 \colon \mathbf{B} \oplus \mathbf{E}\) defined by \(h_0(a) = (a, d_0(a))\) is a unital ring homomorphism (no. 8, Proposition 12); under this mapping, B ⊕ E then becomes an A-algebra. Moreover, if, for all \(x \in \mathbf{M}\) , we write \(h_1(x) = (x, d_1(x))\) , it follows from the definition of \(h_0\) and (29) that \(h_1(ax) = h_0(a)h_1(x)\) ; in other words \(h_1\) is an A-linear mapping of M into B ⊕ E. Then there exists one and only one A-algebra homomorphism, \(h \colon \mathbf{B} \to \mathbf{B} \oplus \mathbf{E}\) such that \(h \mid \mathbf{M} = h_1\) (and necessarily \(h \mid \mathbf{A} = h_0\) ): for, if \(\mathbf{B} = \mathsf{T}(\mathbf{M})\) , this follows from § 5, no. 1, Proposition 1; if \(\mathbf{B} = \mathbb{S}(\mathbf{M})\) , condition (30) shows that \(h(x)h(y) = h(y)h(x)\) for all \(x, y\) in M and the conclusion follows from § 6, no. 1, Proposition 2; finally if \(\mathbf{B} = \bigwedge (\mathbf{M})\) , condition (31) shows that \((h(x))^2 = 0\) for all \(x \in \mathbf{M}\) , since \(x \wedge x = 0\) and the conclusion follows from § 7, no. 1, Proposition 1. The homomorphism \(h\) is such that the composition \(p \circ h \colon \mathbf{B} \to \mathbf{B}\) with the augmentation \(\rho \colon \mathbf{B} \oplus \mathbf{E} \to \mathbf{B}\) is the identity \(l_B\) , for \(p \circ h\) and \(l_B\) coincide by definition for the elements of A and those of M and the set of these elements is a generating system of B. We can therefore write \(h(b) = (b, d(b))\) for all \(b \in B\) and the mapping \(b \mapsto d(b)\) of B into E is a derivation with the required properties, by virtue of Proposition 12 of no. 8.

COROLLARY. Let \(\mathbf{M}\) be a graded \(\mathbf{K}\) -module of type \(\Delta\) ; the \(\mathbf{K}\) -algebras \(\mathsf{T}(\mathbf{M})\) , \(\mathsf{S}(\mathbf{M})\) and \(\bigwedge (\mathbf{M})\) are given the corresponding graduations of type \(\Delta' = \Delta \times \mathbf{Z}\) (§ 5, no. 5, Proposition 7, § 6, no. 6, Proposition 10 and § 7, no. 7, Proposition 11). On the other hand \(\mathbf{M}\) is given the graduation of type \(\Delta'\) such that \(\mathbf{M}_{\alpha,1} = \mathbf{M}_{\alpha}\) for all \(\alpha \in \Delta\) and \(\mathbf{M}_{\alpha,n} = \{0\}\) for \(\alpha \in \Delta\) and \(n \neq 1\) . Let \(\varepsilon'\) be a commutation factor over \(\Delta'\) .

\begin{enumerate}
\def\labelenumi{(\roman{enumi})}
\item
  Let \(\mathbf{E}\) be a graded (left and right) \(\mathsf{T}(\mathbf{M})\) -bimodule of type \(\Delta'\) ; for all \(\delta \in \Delta\) and every integer \(n \in \mathbf{Z}\) , every graded \(\mathbf{K}\) -linear mapping \(f: \mathbf{M} \to \mathbf{E}\) of degree \(\delta_1' = (\delta, n)\) can be extended uniquely to an \(\varepsilon'\) -derivation \(d: \mathsf{T}(\mathbf{M}) \to \mathbf{E}\) of degree \(\delta'\) .
\item
  Let \(\mathbf{E}\) be a graded \(\mathbb{S}(\mathbf{M})\) -module of type \(\Delta'\) ; for a graded \(\mathbf{K}\) -linear mapping \(f: \mathbf{M} \to \mathbf{E}\) of degree \(\delta'\) to be extendable to an \(\varepsilon'\) -derivation \(d: \mathbb{S}(\mathbf{M}) \to \mathbf{E}\) of degree \(\delta'\) , it is necessary and sufficient that, for every ordered pair \((x, y)\) of homogeneous elements of \(\mathbf{M}\) ,
\end{enumerate}

\[
x. f (y) + \varepsilon_ {\delta^ {\prime}, (\deg (y), 1)} ^ {\prime} y. f (x) = y. f (x) + \varepsilon_ {\delta^ {\prime}, (\deg (x), 1)} ^ {\prime} x. f (y).\tag{33}
\]

The \(\varepsilon'\) -derivation \(d\) is then unique.

\begin{enumerate}
\def\labelenumi{(\roman{enumi})}
\setcounter{enumi}{2}
\tightlist
\item
  Let \(\mathbf{E}\) be a graded (left and right) \(\bigwedge (\mathbf{M})\) -bimodule of type \(\Delta'\) ; for a graded K-linear mapping \(f: \mathbf{M} \to \mathbf{E}\) of degree \(\delta'\) to be extendable to an \(\varepsilon'\) -derivation
\end{enumerate}

\[
d: \bigwedge (\mathbf {M}) \rightarrow \mathbf {E}
\]

of degree \(\delta'\) , it is necessary and sufficient that, for every homogeneous element \(x\) of \(\mathbf{M}\) ,

\[
x. f (x) + \varepsilon_ {\delta ^ {\prime}, (\deg (x), 1)} ^ {\prime} f (x). x = 0.\tag{34}
\]

The \(\varepsilon'\) -derivation \(d\) is then unique.

Remark 2 of no. 2 is applied with one of the external B-module laws on E (with B equal to T(M), S(M) or ∧(M)) modified; the external law thus modified is still, by (1) (no. 1), a B-module law and the B-module law thus obtained on E is still compatible with the other B-module structure. It then suffices to apply Proposition 14 with A = K and \(d_{0} = 0\) .

Example (1). In the application of Proposition 14 note that if \(d_0 = 0\) condition (29) means simply that \(d_1\) is A-linear. If we take in particular \(\mathbf{E} = \mathbf{B}\) and the (B, B)-bimodule structure derived from the ring structure on B, conditions (30) and (31) are automatically satisfied when \(d_1\) is taken to be the composition of an endomorphism \(s\) of \(\mathbf{M}\) and the canonical injection \(\mathbf{M} \to \mathbf{B}\) ; this is obvious for (30) since \(\mathbb{S}(\mathbf{M})\) is commutative and for (31) this follows from the fact that \(x\) and \(s(x)\) are of degree 1 in \(\bigwedge (\mathbf{M})\) . It is therefore seen that every endomorphism \(s\) of \(\mathbf{M}\) can be extended uniquely to a derivation \(\mathbf{D}_s\) of \(\mathbb{T}(\mathbf{M})\) (resp. \(\mathbb{S}(\mathbf{M})\) , resp. \(\bigwedge (\mathbf{M})\) ), which is of degree 0. Moreover, for two endomorphisms \(s, t\) of \(\mathbf{M}\) ,

\[
[ \mathbf {D} _ {s}, \mathbf {D} _ {t} ] = \mathbf {D} _ {[ s, t ]}\tag{35}
\]

for both sides are derivations of \(\mathsf{T}(\mathbf{M})\) (resp. \(\mathbb{S}(\mathbf{M})\) , resp. \(\bigwedge (\mathbf{M})\) ) which are equal to \([s, t]\) on \(\mathbf{M}\) .

The expression for \(D_{s}\) is obtained using formula (21) of no. 5, which gives respectively, for \(x_{1}, x_{2}, \ldots, x_{n}\) in M,

\[
\left\{ \begin{array}{l} \mathrm{D} _ {s} (x _ {1} \otimes x _ {2} \otimes \dots \otimes x _ {n}) \\ \qquad = \sum_ {i = 1} ^ {n} x _ {1} \otimes \dots \otimes x _ {i - 1} \otimes s (x _ {i}) \otimes x _ {i + 1} \otimes \dots \otimes x _ {n} \\ \mathrm{D} _ {s} (x _ {1} x _ {2} \dots x _ {n}) = \sum_ {i = 1} ^ {n} x _ {1} \dots x _ {i - 1} s (x _ {i}) x _ {i + 1} \dots x _ {n} \\ \mathrm{D} _ {s} (x _ {1} \wedge x _ {2} \wedge \dots \wedge x _ {n}) \\ \qquad = \sum_ {i = 1} ^ {n} x _ {1} \wedge \dots \wedge x _ {i - 1} \wedge s (x _ {i}) \wedge x _ {i + 1} \wedge \dots \wedge x _ {n}. \end{array} \right.\tag{36}
\]

In the case of the algebra \(\bigwedge (\mathbf{M})\) , there is the following interpretation of \(\mathbf{D}_s\) :

PROPOSITION 15. If M is a free K-module of finite rank n, then, for every endomorphism s of M, the restriction to \(\bigwedge^{n}(\mathbf{M})\) of the derivation \(\mathbf{D}_{s}\) is the homothety of ratio \(\operatorname{Tr}(s)\) .

Let \((e_j)_{1\leqslant j\leqslant n}\) be a basis of M and write \(e = e_1\wedge e_2\wedge \dots \wedge e_n\) . If

\[
s (e _ {j}) = \sum_ {k = 1} ^ {n} \alpha_ {j k} e _ {k},
\]

the third formula in (36) gives

\[
\mathrm{D} _ {s} (e) = \sum_ {i = 1} ^ {n} e _ {1} \wedge \dots \wedge e _ {i - 1} \wedge s \left(e _ {i}\right) \wedge e _ {i + 1} \wedge \dots \wedge e _ {n} = \left(\sum_ {j = 1} ^ {n} \alpha_ {j j}\right) e.
\]

Example (2). In the Corollary to Proposition 14, part (iii), let \(\Delta = \{0\}\) , the graduation on \(\bigwedge (\mathbf{M})\) then being the usual graduation of type \(\mathbf{Z}\) ; on the other hand take \(\varepsilon(p, q) = (-1)^{pq}\) . Then, for every linear form \(x^{*} \in \mathbf{M}^{*}\) on \(\mathbf{M}\) , \(x \mapsto \langle x, x^{*} \rangle\) is a graded K-linear mapping of degree \(-1\) of \(\mathbf{M}\) into \(\bigwedge (\mathbf{M})\) satisfying relation (34); then there exists an antiderivation \(i(x^{*})\) of \(\bigwedge (\mathbf{M})\) , of degree \(-1\) , such that (by virtue of formula (21) of no. 5)

\[
i \left(x ^ {*}\right) \left(x _ {1} \wedge \dots \wedge x _ {n}\right)
\]

\[
= \sum_ {i = 1} ^ {n} (- 1) ^ {i - 1} \left\langle x _ {i}, x ^ {*} \right\rangle x _ {1} \wedge \dots \wedge x _ {i - 1} \wedge x _ {i + 1} \wedge \dots \wedge x _ {n}
\]

and which is a special case of the inner product to be defined in § 11, no. 9, formula (68).

PROPOSITION 16. Let A be a commutative K-algebra, \(\mathbf{M}_i\) ( \(1 \leqslant i \leqslant n\) ) and P A-modules and H the A-module of A-multilinear mappings of \(\mathbf{M}_1 \times \mathbf{M}_2 \times \cdots \times \mathbf{M}_n\) into P. Suppose that there is given a K-derivation \(d_{0}:A\to A\) of the algebra A, for each i, a K-linear mapping \(d_{i}:M_{i}\to M_{i}\) and a K-linear mapping D:P→P, so that, for \(1\leqslant i\leqslant n\) , \((d_{0},d_{i},d_{i})\) is a K-derivation of \((\mathbf{A},\mathbf{M}_{i},\mathbf{M}_{i})\) into itself and \((d_{0},\mathbf{D},\mathbf{D})\) is a K-derivation of \((\mathbf{A},\mathbf{P},\mathbf{P})\) into itself. Then there exists a K-linear mapping \(D^{\prime}:H\to H\) such that \((d_{0},D^{\prime},D^{\prime})\) is a K-derivation of \((\mathbf{A},\mathbf{H},\mathbf{H})\) into itself and

\begin{enumerate}
\def\labelenumi{(\arabic{enumi})}
\setcounter{enumi}{36}
\tightlist
\item
  \(\mathbf{D}(f(x_1, \ldots, x_n))\)
\end{enumerate}

for all \(x_{i} \in \mathbf{M}_{i}\) for \(1 \leqslant i \leqslant n\) and \(f \in \mathbf{H}\) .

We show that for \(f \in H\) , the mapping \(D'f\) of \(M_{1} \times M_{2} \times \cdots \times M_{n}\) into P defined by (37) is A-multilinear. For \(a \in A\) ,

\[
\begin{array}{r l} (\mathrm{D} ^ {\prime} f) (a x _ {1}, x _ {2}, \ldots , x _ {n}) & = \mathrm{D} (a f (x _ {1}, \ldots , x _ {n})) - f (d _ {1} (a x _ {1}), x _ {2}, \ldots , x _ {n}) \\ & \quad - a \sum_ {i = 2} ^ {n} f (x _ {1}, \ldots , x _ {i - 1}, d _ {i} x _ {i}, x _ {i + 1}, \ldots , x _ {n}) \end{array}
\]

and by hypothesis

\[
\mathrm{D} (a f (x _ {1}, \dots , x _ {n})) = (d _ {0} a) f (x _ {1}, \dots , x _ {n}) + a \mathrm{D} (f (x _ {1}, \dots , x _ {n}))
\]

and \(d_{1}(ax_{1}) = (d_{0}a)x_{1} + a.d_{1}x_{1}\) , which gives

\[
(\mathrm{D} ^ {\prime} f) (a x _ {1}, x _ {2}, \dots , x _ {n}) = a. (\mathrm{D} ^ {\prime} f) (x _ {1}, \dots , x _ {n})
\]

and linearity in each of the \(x_{i}\) is proved similarly. On the other hand,

\[
\begin{array}{r l} (\mathrm{D} ^ {\prime} (a f)) (x _ {1}, \dots , x _ {n}) & = \mathrm{D} (a f (x _ {1}, \dots , x _ {n})) \\ & - \sum_ {i = 1} ^ {n} a f (x _ {1}, \dots , x _ {i - 1}, d _ {i} x _ {i}, x _ {i + 1}, \dots , x _ {n}) \\ & = (d _ {0} a) f (x _ {1}, \dots , x _ {n})) + a \mathrm{D} (f (x _ {1}, \dots , x _ {n})) \\ & - \sum_ {i = 1} ^ {n} a f (x _ {1}, \dots , x _ {i - 1}, d _ {i} x _ {i}, x _ {i + 1}, \dots , x _ {n}) \\ & = (d _ {0} a) f (x _ {1}, \dots , x _ {m}) + a (\mathrm{D} ^ {\prime} f) (x _ {1}, \dots , x _ {m}) \end{array}
\]

in other words

\[
\mathbf {D} ^ {\prime} (a f) = \left(d _ {0} a\right) f + a (\mathbf {D} ^ {\prime} f)
\]

which establishes the proposition.

Examples. (3) Applying Proposition 16 to the case \(n = 1\) , \(\mathbf{M}_1 = \mathbf{M}\) , \(P = A\) , then \(H = M^*\) , the dual of \(M\) , and it is seen that for a K-derivation \((d_0, d, d)\) of \((A, M, M)\) is derived a K-derivation \((d_0, d^*, d^*)\) of \((A, M^*, M^*)\) such that

\[
d _ {0} \langle m, m ^ {*} \rangle = \langle d m, m ^ {*} \rangle + \langle m, d ^ {*} m ^ {*} \rangle\tag{38}
\]

for \(m \in \mathbf{M}\) and \(m^* \in \mathbf{M}^*\) . The K-linear mapping of \(\mathbf{M} \oplus \mathbf{M}^*\) into itself which is equal to \(d\) on \(\mathbf{M}\) and to \(d^{*}\) on \(\mathbf{M}^{*}\) then satisfies condition (29) and there is therefore a K-derivation \(\mathbf{D}\) of the A-algebra \(\mathsf{T}(\mathbf{M} \oplus \mathbf{M}^{*})\) , which reduces to \(d_{0}\) on \(\mathbf{A}\) , to \(d\) on \(\mathbf{M}\) and to \(d^{*}\) on \(\mathbf{M}^{*}\) . The restriction \(d_{\mathbf{J}}^{\mathbf{I}}\) of \(\mathbf{D}\) to the sub-A-module \(\mathsf{T}_{\mathbf{J}}^{\mathbf{I}}(\mathbf{M})\) of \(\mathsf{T}(\mathbf{M} \oplus \mathbf{M}^{*})\) ( \(\S 5\) , no. 6) is a K-endomorphism of \(\mathsf{T}_{\mathbf{J}}^{\mathbf{I}}(\mathbf{M})\) such that \((d_{0}, d_{\mathbf{J}}^{\mathbf{I}}, d_{\mathbf{J}}^{\mathbf{I}})\) is a K-derivation of \((\mathbf{A}, \mathsf{T}_{\mathbf{J}}^{\mathbf{I}}(\mathbf{M}), \mathsf{T}_{\mathbf{J}}^{\mathbf{I}}(\mathbf{M}))\) . Moreover, for \(i \in \mathbf{I}\) , \(j \in \mathbf{J}\) , if we write \(\mathbf{I}' = \mathbf{I} - \{i\}\) , \(\mathbf{J}' = \mathbf{J} - \{j\}\) , it is immediately verified that for the contraction \(c_{j}^{t}\) ( \(\S 5\) , no. 6)

\[
c _ {j} ^ {t} (d _ {J} ^ {\mathrm{I}} (z)) = d _ {J ^ {\prime}} ^ {\mathrm{I} ^ {\prime}} (c _ {j} ^ {t} (z)) \quad \text { for   all } z \in \mathbf {T} _ {J} ^ {\mathrm{I}} (\mathbf {M}).
\]

\begin{enumerate}
\def\labelenumi{(\arabic{enumi})}
\setcounter{enumi}{3}
\tightlist
\item
  Let \(M_{i}\) ( \(1 \leqslant i \leqslant 3\) ) be three A-modules and, for each i, suppose that \((d_{0}, d_{i}, d_{i})\) is a derivation of \((\mathbf{A}, \mathbf{M}_{i}, \mathbf{M}_{i})\) ; applying Proposition 16 again for n = 1, a derivation \((d_{0}, d_{ij}, d_{ij})\) of \((\mathbf{A}, \mathbf{H}_{ij}, \mathbf{H}_{ij})\) , where \(\mathrm{H}_{ij} = \mathrm{Hom}_{\mathbf{A}}(\mathbf{M}_{i}, \mathbf{M}_{j})\) , is derived for each ordered pair \((i, j)\) . With this notation, for \(u \in \mathrm{Hom}_{\mathbf{A}}(\mathbf{M}_{1}, \mathbf{M}_{2})\) and \(v \in \mathrm{Hom}_{\mathbf{A}}(\mathbf{M}_{2}, \mathbf{M}_{3})\) ,
\end{enumerate}

\[
d _ {1 3} (v \circ u) = (d _ {2 3} v) \circ u + v \circ (d _ {1 2} u)\tag{39}
\]

as is immediately verified from the definitions.

